% ============================================================
%  EB-1A Research Exhibit Cover Pages
%  Harshavardhan Malla — Cybersecurity Research Portfolio
%  Compile: pdflatex exhibit_cover
% ============================================================

\documentclass[12pt,letterpaper]{article}

\usepackage[margin=1in]{geometry}
\usepackage[T1]{fontenc}
\usepackage[utf8]{inputenc}
\usepackage{lmodern}
\usepackage{booktabs}
\usepackage{array}
\usepackage{xcolor}
\usepackage{fancyhdr}
\usepackage{titlesec}
\usepackage{enumitem}
\usepackage{hyperref}
\usepackage{mdframed}
\usepackage{amsmath}
\usepackage{microtype}
\usepackage{graphicx}

\definecolor{headblue}{RGB}{0,51,102}
\definecolor{lightblue}{RGB}{230,240,255}
\definecolor{darkgray}{RGB}{80,80,80}

\pagestyle{fancy}
\fancyhf{}
\renewcommand{\headrulewidth}{0.5pt}
\fancyhead[L]{\small\textcolor{darkgray}{EB-1A Research Portfolio --- Harshavardhan Malla}}
\fancyhead[R]{\small\textcolor{darkgray}{\thepage}}

\titleformat{\section}{\large\bfseries\color{headblue}}{}{0em}{}[\vspace{-0.5em}\hrule\vspace{0.5em}]

\newenvironment{exhibitbox}{%
  \begin{mdframed}[backgroundcolor=lightblue,linecolor=headblue,linewidth=1.5pt,
                   topline=true,bottomline=true,leftline=true,rightline=true,
                   innertopmargin=12pt,innerbottommargin=12pt,
                   innerleftmargin=14pt,innerrightmargin=14pt]
}{%
  \end{mdframed}
}

\newcommand{\exhibitlabel}[1]{%
  \begin{center}
    \large\bfseries\textcolor{headblue}{#1}
  \end{center}
}

% ============================================================
\begin{document}

\begin{titlepage}
  \centering
  \vspace*{2cm}

  {\Huge\bfseries\textcolor{headblue}{Research Portfolio}}\\[0.5em]
  {\Large Evidence of Extraordinary Ability in Cybersecurity Research}\\[2em]

  \hrule height 2pt
  \vspace{1em}

  {\LARGE\bfseries Harshavardhan Malla}\\[0.5em]
  {\large Independent Researcher}\\[0.5em]
  {\large Cybersecurity --- Vulnerability Management \& Anomaly Detection}\\[2em]

  \hrule height 2pt
  \vspace{2em}

  \begin{tabular}{ll}
    \textbf{Portfolio date:}   & May 31, 2026 \\[0.3em]
    \textbf{Papers included:}  & 4 original research papers \\[0.3em]
    \textbf{Venues targeted:}  & IEEE, USENIX, ACM (peer-reviewed) \\[0.3em]
    \textbf{Data policy:}      & All data synthetic and seeded; \\
                                & no employer or production data used \\[0.3em]
    \textbf{Code policy:}      & All code and results publicly released \\
  \end{tabular}

  \vspace{3em}

  \begin{exhibitbox}
    \textbf{Contents of This Portfolio:}
    \begin{enumerate}[leftmargin=*]
      \item \textbf{Exhibit A} --- Paper 1: VulnPrio --- Vulnerability Prioritization Benchmark
      \item \textbf{Exhibit B} --- Paper 2: CalibScore --- Failure-Aware Calibration Gate
      \item \textbf{Exhibit C} --- Paper 3: HygieneBench --- Cyber-Hygiene Anomaly Detection Benchmark
      \item \textbf{Exhibit D} --- Paper 4: HygienePrio --- Hygiene-Augmented Exploit-Weighted Scorer
    \end{enumerate}

    \vspace{0.5em}
    Each exhibit includes: (1) a one-page layperson significance statement;
    (2) a technical abstract; and (3) the full paper draft.
  \end{exhibitbox}

\end{titlepage}

% ============================================================
% EXHIBIT A COVER
% ============================================================
\newpage
\setcounter{page}{1}

\exhibitlabel{EXHIBIT A}
\section{Paper 1: VulnPrio}

\begin{center}
  \textbf{\large Context-Aware Vulnerability Prioritization for Government Endpoint Fleets:\\
  Integrating Exploit Intelligence, Asset Criticality, and Endpoint Telemetry}\\[0.5em]
  \textit{Harshavardhan Malla} --- Independent Researcher\\[0.5em]
  Target venue: IEEE/USENIX Security Workshop on Measurement and Practice
\end{center}

\vspace{1em}
\noindent\textbf{Layperson Summary:}

Government agencies receive thousands of new software vulnerability reports each month
but can only fix a fraction of them before the next wave arrives. The standard approach
--- ranking by ``severity score'' (CVSS) --- does not account for whether a vulnerability
is actually being exploited in the wild, how critical the affected computer is, or how
complex the patch is to apply. This paper introduces \textbf{VulnPrio}, a mathematical
framework that combines seven signals (exploit probability, known exploitation status,
severity, asset criticality, network exposure, urgency, and remediation difficulty)
into a single priority score, and embeds that score in a scheduling system that
automatically handles regulatory deadlines and records every decision in a
tamper-evident log.

The key finding is deliberately honest: under the current research conditions,
VulnPrio performs no better than simpler approaches --- and the paper says so
explicitly. This ``null result'' reported with full statistical rigor is the
contribution: it establishes a reproducible testing framework that any future
system can be compared against, and proves that data-driven superiority claims
require real calibrated data --- not placeholder weights.

\vspace{0.5em}
\noindent\textbf{Field Impact:}
\begin{itemize}[leftmargin=*,itemsep=2pt]
  \item First benchmark combining EPSS signals, capacity constraints,
        KEV-deadline overrides, and hash-chained audit logs in one reproducible pipeline
  \item 390 hash-chained audit logs, all verifying without error ---
        a tamper-evident decision record mechanism suitable for regulatory environments
  \item Honest null-result reporting with full statistical apparatus
        (Wilcoxon + Holm + BCa bootstrap) across 30 seeds and 13 strategies
  \item Open code and frozen results released for community replication
\end{itemize}

\vspace{0.5em}
\noindent\textbf{Technical Abstract:}

\begin{exhibitbox}
Vulnerability remediation in large government endpoint fleets is constrained by
finite patch capacity, mandatory regulatory timelines, and the daily flood of
newly published CVEs. We present \textbf{VulnPrio}, a seven-feature linear
prioritization framework integrating EPSS exploit probability, KEV status, CVSS
base score, asset criticality, network exposure, urgency, and remediation complexity.
The framework is embedded in a five-phase capacity-constrained scheduler with
KEV-deadline overrides and append-only SHA-256 hash-chained audit logs. Evaluated
against 12 competing strategies across 30 seeds on a fully synthetic fleet dataset,
all 13 strategies are \emph{statistically indistinguishable} on the Expected
Exploited Host-Days (EHD) metric (inter-strategy spread $<$0.5\%; all 78 pairwise
Wilcoxon+Holm tests: adjusted $p>0.05$). We report this null result honestly:
the contribution is a reproducible, falsifiable, audit-evidence-producing benchmark,
not a performance claim. All 390 audit logs hash-verify correctly.
\end{exhibitbox}

% ============================================================
% EXHIBIT B COVER
% ============================================================
\newpage

\exhibitlabel{EXHIBIT B}
\section{Paper 2: CalibScore}

\begin{center}
  \textbf{\large When Calibration Fails: A Failure-Aware Public-Feed Gate for\\
  Vulnerability Prioritization Under Sparse Exploit Labels}\\[0.5em]
  \textit{Harshavardhan Malla} --- Independent Researcher\\[0.5em]
  Target venues: USENIX CSET (primary); LASER Workshop; ACM DTRAP
\end{center}

\vspace{1em}
\noindent\textbf{Layperson Summary:}

Paper 1 uses placeholder (estimated) weights in its scoring formula because no
one had measured whether real public data could calibrate them. Paper 2 asks that
question directly: can we use publicly available exploit records to find better
weights for a realistic government fleet? The answer is no --- and this paper
explains precisely why, with a formal measurement.

The paper introduces a \textbf{failure-aware gate}: a step-by-step procedure
for deciding whether you have enough data to calibrate a model before you try.
The gate was written down (pre-registered) before any data was collected, so
there was no possibility of changing the rules after seeing the results. When
applied to a realistic 31-product government fleet, the gate found only 7
confirmed exploitation events over two years --- far below the minimum of 20
needed to safely fit a calibration model. The gate said ``PIVOT --- do not
calibrate,'' and the paper reports exactly that outcome as its primary finding.

\vspace{0.5em}
\noindent\textbf{Field Impact:}
\begin{itemize}[leftmargin=*,itemsep=2pt]
  \item First paper to propose a formal ``feasibility gate'' before attempting
        public-feed calibration of vulnerability prioritization weights
  \item Pre-registration of analysis plan (8 stop rules) before any API queries
        were executed --- a clinical-trial-grade methodological standard applied to
        cybersecurity research
  \item Quantified data-density boundary: provides a scaling formula
        ($E[\text{positives}] \approx p \cdot r \cdot m \cdot T$) for practitioners
        to estimate feasibility before committing to a calibration study
  \item BCa confidence intervals so wide (Top-10 coverage: [0.00, 0.71]) that
        they honestly communicate irreducible uncertainty --- no false precision
\end{itemize}

\vspace{0.5em}
\noindent\textbf{Technical Abstract:}

\begin{exhibitbox}
Exploit-occurrence labels from the CISA KEV catalog are sparse in bounded product
catalogs. We contribute a pre-registered, multi-stage gate methodology for
assessing whether public-feed signals (EPSS v3, NVD, KEV) are dense enough to
calibrate vulnerability-prioritization weights without label leakage. The gate
comprises: chunked API acquisition with checkpointing; an 18-window multi-$t_0$
aggregation that maximizes observable distinct positive CVEs; and a three-tier
decision rule (PROCEED $\geq$50, CAUTION 20--49, PIVOT $<$20) backed by a
pre-registered stop-rule registry. Applied to 2,688 catalog-matched CVEs over
the EPSS v3 era (2023-03-07 to 2025-03-16), we observe only \textbf{7 unique
positive CVEs} across 18 windows --- triggering a gate decision of
\textbf{PIVOT}: calibration is not attempted. The gate PIVOT is the primary
contribution. A pre-registered sensitivity sweep of six design-prior weight
vectors produces descriptive statistics with wide BCa CIs, confirming that
no vector can be distinguished under this sparsity regime.
\end{exhibitbox}

% ============================================================
% EXHIBIT C COVER
% ============================================================
\newpage

\exhibitlabel{EXHIBIT C}
\section{Paper 3: HygieneBench}

\begin{center}
  \textbf{\large HygieneBench: A Reproducible Synthetic Benchmark for\\
  Cyber-Hygiene Anomaly Detection Across Identity, Endpoint, and Patch Telemetry}\\[0.5em]
  \textit{Harshavardhan Malla} --- Independent Researcher\\[0.5em]
  Target venues: ACM CCS AISec Workshop (primary); IEEE S\&P Workshop
\end{center}

\vspace{1em}
\noindent\textbf{Layperson Summary:}

Good cyber security depends on basic hygiene: making sure old accounts are
closed, computer agents are running, patches are applied on time, and data
is being collected reliably. When these hygiene checks fail, the conditions
exist for attackers to exploit dormant accounts, unmonitored computers, or
unpatched vulnerabilities. Detecting hygiene failures automatically is
therefore a fundamental security task --- but there is no public testing
dataset for this problem. Researchers developing AI tools for hygiene
monitoring have had no standard way to compare their methods.

\textbf{HygieneBench} fills this gap. It generates realistic (but fully
synthetic) organizational data spanning employee accounts, computers, software
patches, vulnerability records, and data freshness logs. It includes 12 types
of hygiene anomalies and 7 detection tasks, and it tests 8 different detection
methods across 5 conditions (including ``stale data'' conditions that simulate
real-world data pipeline failures). The central finding is surprising:
\textbf{simple rule-based detection} outperforms sophisticated machine learning
on 86.2\% of tested configurations. Only two tasks (detecting unusual account
group changes, and detecting clusters of unpatched critical vulnerabilities)
consistently benefit from machine learning.

\vspace{0.5em}
\noindent\textbf{Field Impact:}
\begin{itemize}[leftmargin=*,itemsep=2pt]
  \item \textbf{First public benchmark} jointly covering Active Directory identity
        hygiene, endpoint patch posture, vulnerability exposure, and telemetry
        freshness as first-class evaluation dimensions
  \item Fully open, seeded synthetic generator grounded in public structural priors
        (NIST NVD, Verizon DBIR 2026, CISA BOD 23-01) ---
        reproducible across platforms and research groups
  \item Pre-registered failure protocol ($\Delta \geq 0.05$ AP in $\geq$2/3 seeds) applied
        rigorously across all 810 evaluation runs --- 86.2\% failure rate reported
        without filtering or selective presentation
  \item Telemetry staleness quantified: $-0.168$ AP on endpoint--identity correlation
        under 20\% stale entity rate --- the first benchmark-level evidence for
        the detection cost of stale data pipelines
  \item Motivates three changes to benchmark practice: task-stratified reporting,
        negative-result protocols, and telemetry freshness as a mandatory axis
\end{itemize}

\vspace{0.5em}
\noindent\textbf{Technical Abstract:}

\begin{exhibitbox}
Cyber-hygiene anomaly detection is operationally important but poorly served by
existing benchmarks. We introduce \textbf{HygieneBench}, an open, reproducible
benchmark covering seven evaluation tasks and twelve anomaly classes across
eleven entity/event tables of synthetic Active Directory, endpoint, and
vulnerability telemetry. All data are generated from a seeded synthetic pipeline
grounded in public structural priors (NIST NVD, Verizon DBIR 2026, CISA BOD 23-01).
An 810-run evaluation of eight detection methods under five telemetry conditions,
using a pre-registered failure protocol ($\Delta \geq 0.05$ AP in $\geq$2/3 seeds),
finds that \textbf{86.2\% of configurations fail to consistently outperform the
rule baseline}. ML adds meaningful signal on two of seven tasks: group-membership
drift (T2, +0.185 AP, temporal z-score) and patch-vulnerability hygiene
(T5, +0.210 AP, OCSVM). Telemetry staleness (C-STALE) consistently degrades
performance ($-$0.168 AP on T3). Generator, datasets, and harness are openly
released.
\end{exhibitbox}

% ============================================================
% EXHIBIT D COVER
% ============================================================
\newpage

\exhibitlabel{EXHIBIT D}
\section{Paper 4: HygienePrio}

\begin{center}
  \textbf{\large HygienePrio: Cyber-Hygiene Signal Augmentation for\\
  EPSS-Weighted Vulnerability Prioritization}\\[0.5em]
  \textit{Harshavardhan Malla} --- Independent Researcher\\[0.5em]
  Target venues: IEEE Transactions on Network and Service Management (TNSM);
  ACM Digital Threats: Research and Practice (DTRAP)
\end{center}

\vspace{1em}
\noindent\textbf{Layperson Summary:}

Paper 1 showed that combining public vulnerability signals (exploit likelihood,
CVSS severity, KEV status) did not improve over a simple exploit-likelihood
baseline when only CVE-level features were used. Paper 3 showed that host-level
hygiene signals --- how many patches are overdue, how exposed the machine is
through network accounts, how recently data was collected from it --- carry
genuine discriminative structure in a realistic synthetic fleet. \textbf{HygienePrio}
connects these two findings: it asks whether the ``missing signal'' from Paper 1
was host-level hygiene posture rather than more CVE-level detail.

The answer, in the synthetic evaluation, is yes. By building a three-dimensional
Hygiene Risk Score (patch debt, Active Directory exposure, telemetry freshness)
and combining it with exploit likelihood through a principled weighted formula,
HygienePrio achieves substantially higher Precision at the top-50 priority slots
than EPSS alone (50.9\% vs.\ 20.2\% in the synthetic evaluation; +31 percentage
points). Removing the patch-posture component costs the most ground (-20 pp),
confirming Paper 3's finding that patch-vulnerability hygiene is the most
discriminative hygiene signal. All results are fully qualified as synthetic-fleet
findings; real-fleet validation is explicitly identified as future work.

\vspace{0.5em}
\noindent\textbf{Field Impact:}
\begin{itemize}[leftmargin=*,itemsep=2pt]
  \item \textbf{First scorer} to jointly optimize per-CVE exploit likelihood
        (EPSS) and host-level hygiene posture (patch debt, AD exposure,
        telemetry freshness) for (host, CVE) pair prioritization
  \item Closes the null-result loop opened by Paper 1: demonstrates that
        host-level hygiene signals --- not additional CVE-level features ---
        are the missing link for meaningful prioritization improvement
  \item Pre-registered 25-seed evaluation with BCa 95\% CIs; no post-hoc
        metric selection; ground-truth definition fixed before scoring
  \item Ground-truth circularity (HRS in both scorer and labels)
        explicitly acknowledged in both the pre-registration and the
        threats-to-validity section --- clinical-trial-grade transparency
  \item Open Python implementation (HygienePrio scorer + EEHDA fleet generator
        + 225-row frozen results) released for community replication
\end{itemize}

\vspace{0.5em}
\noindent\textbf{Technical Abstract:}

\begin{exhibitbox}
Capacity-constrained vulnerability remediation teams must prioritize (host, CVE)
pairs for patching within each maintenance window. The Exploit Prediction Scoring
System (EPSS) is asset-agnostic by design: it scores a CVE, not a host. Prior
work (\textbf{VulnPrio}, Paper 1) found that augmenting EPSS with CVE-level
features did not improve prioritization, suggesting the missing signal is
host-level hygiene posture. We present \textbf{HygienePrio}, a scoring framework
integrating a three-dimensional Hygiene Risk Score (HRS) --- patch posture, Active
Directory exposure, and telemetry freshness --- into an EPSS-weighted scorer via
an additive combination with an explicit EPSS$\times$HRS interaction term.
In a pre-registered synthetic evaluation across 25 fleet seeds of the EEHDA
synthetic fleet, HygienePrio achieves mean Precision@50 of 0.509 compared with
EPSS-only's 0.202 (approximately +31 pp; non-overlapping BCa 95\% CIs; synthetic
evaluation only). Leave-one-dimension-out ablation confirms patch posture as the
dominant HRS component ($-$20 pp when removed). All results are bounded to the
synthetic evaluation context; generalization to real enterprise fleets requires
external validation not performed here.
\end{exhibitbox}

% ============================================================
% RESEARCH PROGRAM OVERVIEW
% ============================================================
\newpage

\section{Research Program Coherence}

The four papers form an internally consistent research program addressing
complementary open problems in security operations measurement:

\vspace{1em}
\begin{center}
\begin{tabular}{>{\bfseries\raggedright\arraybackslash}p{3.5cm}
                >{\raggedright\arraybackslash}p{4.2cm}
                >{\raggedright\arraybackslash}p{4.0cm}
                >{\raggedright\arraybackslash}p{2.8cm}}
\toprule
\textbf{Paper} & \textbf{Open Problem} & \textbf{Contribution} & \textbf{Finding} \\
\midrule
Paper 1 \newline (VulnPrio) &
  How do we prioritize patches in a capacity-constrained, audit-required environment? &
  7-feature composite score + 5-phase scheduler + hash-chained audit log + 30-seed reproducible benchmark &
  Null: all 13 strategies indistinguishable under placeholder weights \\
\addlinespace
Paper 2 \newline (CalibScore) &
  Can real public-feed labels calibrate the Paper 1 weights? &
  Pre-registered failure-aware gate + multi-t\textsubscript{0} aggregation + stop-rule registry &
  Negative: 7 unique positives, gate PIVOT, calibration not attempted \\
\addlinespace
Paper 3 \newline (HygieneBench) &
  Do ML methods improve over rule baselines for hygiene anomaly detection? &
  First open hygiene benchmark + 12-class taxonomy + 7-task suite + 5-condition evaluation &
  Mixed: ML adds value on 2/7 tasks; rules dominate on 5/7 tasks \\
\addlinespace
Paper 4 \newline (HygienePrio) &
  Do host-level hygiene signals (patch posture, AD exposure, telemetry freshness) improve EPSS-weighted (host, CVE) prioritization? &
  HygienePrio scorer with HRS$\times$EPSS interaction term + pre-registered 25-seed evaluation + open Python implementation &
  Positive (synthetic): $+$31 pp P@50 over EPSS-only; patch posture dominant (noPatch: $-$20 pp) \\
\bottomrule
\end{tabular}
\end{center}

\vspace{1.5em}
\noindent\textbf{Common methodological principles across all four papers:}

\begin{enumerate}[leftmargin=*,itemsep=4pt]
  \item \textbf{Pre-registration:} Analysis plans committed before data examination
        (Paper 2: 8 stop rules; Paper 3: failure protocol, k values, feature sets;
        Paper 4: HRS weights, calibration grid, evaluation seeds, ground-truth definition)
  \item \textbf{Reproducibility:} All results derivable from open code + seeded generator;
        no result depends on closed tools or private data
  \item \textbf{Honest reporting:} Null and negative results reported as primary
        contributions, not buried as limitations
  \item \textbf{Statistical rigor:} Wilcoxon + Holm correction for multiple comparisons;
        BCa bootstrap CIs; paired-seed design to control seed variance
  \item \textbf{Ethical data use:} All data fully synthetic and seeded;
        no employer, government, or production data used at any stage
\end{enumerate}

\vspace{1.5em}
\noindent\textbf{Evidence of original contribution:}

Each paper identifies a genuine gap in the security research literature
(acknowledged in Related Work with specific citations), proposes a concrete
methodology to address it, executes the methodology with full reproducibility,
and reports findings honestly regardless of direction. The research demonstrates:

\begin{itemize}[leftmargin=*,itemsep=4pt]
  \item Ability to formulate novel research questions in specialized domains
  \item Mastery of experimental design appropriate to the specific question
  \item Disciplined statistical practice applied to security measurement problems
  \item Infrastructure contributions (benchmarks, datasets, pipelines) enabling community research
  \item Intellectual integrity through transparent null and negative result reporting
\end{itemize}

\end{document}
